The main implication is that epoch-aware thresholding provides a clear advantage over naive baselines for blink detection in epoch-structured EEG data. Proposed-Med achieved a macro-$F_1$ of 0.8392, compared with 0.7242 for BLINKER-concat and 0.5731 for MNE-annot, indicating that adapting the decision rule to the epoch structure improved performance over applying baseline detectors without such adaptation. This pattern is consistent with the observation that running a continuous-signal detector such as BLINKER on epoch-structured data underperformed epoch-aware estimation \citep{kleifges2017blinker}. At the same time, both proposed configurations significantly outperformed BLINKER-concat and MNE-annot under Wilcoxon tests with Bonferroni correction, with Proposed-Med obtaining the highest macro-$F_1$ overall, although its margin over Proposed-Mean was small and not statistically significant.
